\chapter{Proposed Scheme 研究方法與系統設計}

\section{問題定義與需求結構}
\label{sec:problem_definition}

系統輸入為使用者的自然語言需求、出發位置與時間，輸出為最多三間不同餐廳，以及各自的情境依據與建議公車方案。最終決策對象是餐廳；同一餐廳可能有不同路線，系統先找出合適路線，再以交通條件回饋餐廳排序。

需求分為語意偏好、必要限制與動態狀態。「適合約會」「慶生」用於評論語意召回；預算、步行上限與最多可轉乘次數作為必要條件；公車 ETA、抵達時間與營業狀態則隨查詢時間分析。相對時間須轉為具體日期與時刻，預算也須區分每人與整團；關鍵欄位不明時先補問。

為分析交通，暫以 \(z=(r,p)\) 表示餐廳 \(r\) 與公車方案 \(p\) 的組合。這個表示用於計算與查核，不代表餐廳與交通是兩個同等的推薦目的。外部文字僅作推薦證據，不具有修改系統規則或觸發額外工具的權限。


需求以
\[
q=(c,o,t_{\mathrm{dep}},B,L_{\max},W_{\max},N_{\max},D_{\mathrm{meal}})
\]
表示，其中 \(c\) 為用餐情境，\(o\) 為出發位置，\(B\) 為同單位預算，
\(L_{\max}\) 為步行上限，\(W_{\max}\) 為候車上限，\(N_{\max}\) 為轉乘上限，
\(D_{\mathrm{meal}}\) 為用餐時長；步行上限另依使用者指定的單段或總量定義保存。
系統先由文字語料 \(\mathcal{D}\) 召回餐廳集合
\(\mathcal{R}_{K}(q)\)，再為每間餐廳取得路線集合
\(\mathcal{P}(r,o,t_{\mathrm{dep}})\)。
因此，語意召回決定搜尋範圍，交通分析決定此範圍內各餐廳的可行路線與最終次序。

{\small\singlespacing
\renewcommand{\arraystretch}{1.25}
\setlength{\tabcolsep}{3pt}
\begin{longtable}{|P{\dimexpr 0.1700\linewidth-2\tabcolsep-1.3333\arrayrulewidth\relax}|P{\dimexpr 0.4000\linewidth-2\tabcolsep-1.3333\arrayrulewidth\relax}|P{\dimexpr 0.4300\linewidth-2\tabcolsep-1.3333\arrayrulewidth\relax}|}
\caption{需求條件與處理分工}\label{table:condition_types}\\
\hline
\textbf{類型} & \textbf{表示內容} & \textbf{系統角色} \\ \hline
\endfirsthead
\hline
\textbf{類型} & \textbf{表示內容} & \textbf{系統角色} \\ \hline
\endhead
情境偏好 & 約會、慶生、朋友聚餐及評論片段 & 語意召回與相關性排序 \\ \hline
必要條件 & 同單位預算、營業區間、步行與轉乘上限 & 以證據逐項判為 T、F 或 U \\ \hline
交通狀態 & ETA、候車、轉乘及抵達時間 & 建立時間線並回饋餐廳篩選 \\ \hline
\end{longtable}
}


\section{設計原則}

第一，先找情境相符的餐廳，再安排交通，避免因為交通方便而推薦明顯不符情境的餐廳。第二，在最終排名前分析交通，避免只在既定餐廳後附一條路線。第三，必要限制由程式查核，較高語意分數不能抵銷明確不符合的預算或搭乘條件。第四，來源不足時表達不確定，不把資料缺漏當成已確認可行。第五，推薦說明須能回到評論、資料與判定紀錄。

初版以查詢當下的餐廳推薦與交通比較為核心，行程中監控、延誤重規劃與回程檢查列於 \ref{sec:extensions} 節的後續擴充。

\section{系統架構}
\label{sec:architecture}

圖~\ref{fig:architecture} 由九個編號模組組成。模組 1--3 為離線資料整備，模組 4--9 為線上推薦流程。


% 架構與流程圖共用配色；在 figure 群組外載入。
\usetikzlibrary{arrows.meta,shapes.geometric}
\definecolor{archLLM}{HTML}{DCE9F8}
\definecolor{archLLMd}{HTML}{2F6DB5}
\definecolor{archPro}{HTML}{E6F2E6}
\definecolor{archProd}{HTML}{3C8D4F}
\definecolor{archKB}{HTML}{FFF3D6}
\definecolor{archKBd}{HTML}{C58A12}
\definecolor{archSrc}{HTML}{FDE9DC}
\definecolor{archSrcd}{HTML}{D9622B}
\definecolor{archGray}{HTML}{4A4A4A}

\begin{figure}[H]
\singlespacing
\centering
\resizebox{\textwidth}{!}{% 主流程由左至右；離線資料與外部服務由上方提供輸入。
% 架構與流程圖共用配色；在 figure 群組外載入。
\usetikzlibrary{arrows.meta,shapes.geometric}
\definecolor{archLLM}{HTML}{DCE9F8}
\definecolor{archLLMd}{HTML}{2F6DB5}
\definecolor{archPro}{HTML}{E6F2E6}
\definecolor{archProd}{HTML}{3C8D4F}
\definecolor{archKB}{HTML}{FFF3D6}
\definecolor{archKBd}{HTML}{C58A12}
\definecolor{archSrc}{HTML}{FDE9DC}
\definecolor{archSrcd}{HTML}{D9622B}
\definecolor{archGray}{HTML}{4A4A4A}

\begin{tikzpicture}[font=\normalsize,>=Stealth,
  box/.style={draw=archProd,fill=archPro,line width=0.8pt,rounded corners=4pt,
    minimum width=3.0cm,minimum height=2.1cm,text width=2.8cm,align=center,inner sep=5pt},
  llm/.style={box,draw=archLLMd,fill=archLLM},
  kb/.style={box,draw=archKBd,fill=archKB},
  src/.style={box,draw=archSrcd,fill=archSrc,dashed,minimum height=1.3cm},
  badge/.style={circle,fill=archGray,text=white,font=\normalsize\bfseries,inner sep=2pt,minimum size=0.55cm},
  arr/.style={-{Stealth[length=2.4mm,width=2mm]},line width=1pt,archGray},
  fb/.style={arr,dashed,archSrcd}]
\fill[archPro!35,draw=archProd,dashed,rounded corners=7pt]
  (-1.75,4.9) rectangle (8.75,8.1);
\node[anchor=north west,text=archProd,font=\normalsize\bfseries] at (-1.5,8.0) {離線資料整備};
\node[box] (n1) at (0,6.3) {\textbf{資料蒐集}\\評論、影片\\與店家資訊};
\node[box] (n2) at (3.5,6.3) {\textbf{屬性與證據整理}\\實體、情境\\價格與時間};
\node[kb] (n3) at (7,6.3) {\textbf{資料儲存}\\圖譜、關聯表\\與向量索引};
\node[src] (u) at (0,3.15) {\textbf{使用者需求}\\情境、位置與時間};
\node[src] (tdx) at (10.5,3.15) {\textbf{TDX 資料}\\旅運規劃與公車};
\node[anchor=west,text=archGray,font=\normalsize\bfseries,align=left] at (12.6,3.15) {線上推薦\\餐廳為最終決策單位};
\node[llm] (n4) at (0,0) {\textbf{需求理解與規劃器}\\LLM 解析\\程式安排工具};
\node[box] (n5) at (3.5,0) {\textbf{評論語意召回}\\情境相符的\\候選餐廳};
\node[box] (n6) at (7,0) {\textbf{領域證據與粗篩}\\價格、營業\\與必要條件};
\node[box] (n7) at (10.5,0) {\textbf{公車交通分析}\\步行、候車\\乘車與轉乘};
\node[box] (n8) at (14,0) {\textbf{餐廳篩選與排序}\\三值查核\\情境與交通};
\node[llm] (n9) at (17.5,0) {\textbf{輸出與互動}\\餐廳、評論依據\\交通與未知項};
\foreach \i in {1,...,9}{\node[badge] at (n\i.north west) {\i};}
\draw[arr] (n1)--(n2);
\draw[arr] (n2)--(n3);
\draw[arr] (u)--(n4);
\foreach \a/\b in {4/5,5/6,6/7,7/8,8/9}{\draw[arr] (n\a.east)--(n\b.west);}
\draw[arr] (n3.south)--(7,4.45)--(3.5,4.45)--(n5.north);
\draw[arr] (n3.south)--(n6.north);
\draw[arr,archSrcd] (tdx.south)--(n7.north);
\draw[fb] (n9.south)--(17.5,-2.0)--(0,-2.0)--(n4.south);
\node[fill=white,inner sep=3pt,text=archSrcd,font=\small] at (8.75,-2.0) {使用者修改需求後重新查詢};
\begin{scope}[shift={(0,-3.3)}]
 \draw[archLLMd,fill=archLLM] (0,0) rectangle (0.45,0.3);
 \node[anchor=west,font=\small] at (0.6,0.15) {語言模型與互動};
 \draw[archProd,fill=archPro] (5.3,0) rectangle (5.75,0.3);
 \node[anchor=west,font=\small] at (5.9,0.15) {程式分析};
 \draw[archKBd,fill=archKB] (9.5,0) rectangle (9.95,0.3);
 \node[anchor=west,font=\small] at (10.1,0.15) {資料儲存};
\end{scope}
\end{tikzpicture}
}
\caption[系統架構圖]{系統架構圖：模組 1--3 建立餐廳與證據資料，模組 4--9 由左至右依序完成需求理解、評論語意召回、公車交通分析與餐廳推薦。交通條件在最終排名前納入，行程中監控屬後續擴充。}
\label{fig:architecture}
\end{figure}


模組 1 蒐集可用的餐廳評論、影片逐字稿與店家資訊。模組 2 整理餐廳實體、分店、情境片段、價格與營業資料，保存來源並人工核對。模組 3 將文字片段存入向量索引，將餐廳、屬性與證據關係存入圖譜及關聯表。

模組 4 由語言模型將需求轉為固定欄位；程式化規劃器檢查輸入、安排工具並處理失敗。模組 5 以關鍵字、向量或混合方法對齊情境與評論，按餐廳去重形成候選集合。
混合召回可採正規化分數的加權融合：
\[
s_{\mathrm{retrieve}}(r,q)=\lambda\widetilde{s}_{\mathrm{BM25}}(r,q)
+(1-\lambda)\widetilde{s}_{\mathrm{vector}}(r,q),\qquad 0\leq\lambda\leq1.
\]
其中兩類分數均先按固定規則正規化，文字片段依餐廳分組後再形成候選；
召回數量 \(K\)、最低相關性門檻、正規化方式、片段聚合規則與 \(\lambda\) 必須版本化，避免評論數量較多的餐廳僅因重複片段占據候選。
模組 6 查詢相關證據，先排除已知預算等非路線條件不符者；完整營業與步行判定留待交通資料取得後處理。

模組 7 以出發地、時間及餐廳位置取得公車候選，分析步行、候車、乘車與轉乘時間。模組 8 對餐廳與路線組合查核，排除明確不符者，未知者另列；為每間餐廳保留合適方案，再結合語意與交通排序。模組 9 顯示餐廳、評論語意依據、交通負擔與來源；使用者可修正情境或交通條件，重新進入相同流程。

\section{資料來源與品質控管}
\label{sec:data}


{\small\singlespacing
\renewcommand{\arraystretch}{1.25}
\setlength{\tabcolsep}{3pt}
\begin{longtable}{|P{\dimexpr 0.1900\linewidth-2\tabcolsep-1.3333\arrayrulewidth\relax}|P{\dimexpr 0.3200\linewidth-2\tabcolsep-1.3333\arrayrulewidth\relax}|P{\dimexpr 0.4900\linewidth-2\tabcolsep-1.3333\arrayrulewidth\relax}|}
\caption{資料來源與用途}\label{table:data_sources}\\
\hline
\textbf{資料} & \textbf{預計來源} & \textbf{用途與核對方式} \\ \hline
\endfirsthead
\hline
\textbf{資料} & \textbf{預計來源} & \textbf{用途與核對方式} \\ \hline
\endhead
餐廳評論與情境描述 & 可合法取得的評論資料；原稿規劃的 YouTube 逐字稿與留言可作補充 & 以文字片段對齊約會、慶生等情境；保存餐廳 ID、來源與片段，不把主觀描述當成事實保證 \\ \hline
餐廳名稱與分店 & 店家公開資訊、人工核對；Places API 協助地點對應 & 同名或多分店時保留核對，不只靠名稱合併 \\ \hline
營業時段與價格 & 店家官方資訊及必要的人工核對 & 保存核對日期、適用日期與原始價格區間 \\ \hline
餐廳座標 & 可追溯的位置資料或現場核對 & 用於路線查詢；Place ID 不能代替座標 \\ \hline
公車搭乘方案 & TDX 旅運規劃 & 取得各候選餐廳的公車與步行路線，再分析條件 \\ \hline
公車 ETA 與狀態 & TDX 預估到站 N1 & 查核當下可取得的上車與轉乘依據 \\ \hline
站位、路線、班表與站間時間 & TDX 公車資料 & 核對方向、站序、服務日及未來服務，區分班表與即時依據 \\ \hline
\end{longtable}
}


評論是本研究的核心語意資料。影片逐字稿與留言規劃為補充來源；正式評論來源與可使用範圍須先確認，不能把影片敘述、觀眾留言與到店評論混作同一種證據。資料須標明文字類型，保留各自的來源與用途。

每筆文字與事實保存來源 ID、可得的發布時間、觀測或更新時間、本系統取得時間及適用日期。缺少時間欄位時保留空值，不以抓取時間冒充來源更新時間；影片發布日期也不能證明店家價格至今未變。

本研究僅規劃以 Places API 進行地點對應，並持久保存 Place ID；不預設可以大量保存其評論建立向量語料。若評論來源改用平台 API，須先核對適用的保存、索引與展示政策，再確認設計~\cite{google_places_policy}。

每項判定另存 \texttt{restaurant\_id}、\texttt{plan\_id}、\texttt{constraint\_id}、T／F／U、reason、\texttt{evidence\_ids} 與判定時點。圖譜支援餐廳、情境屬性及來源片段的查詢，向量索引負責語意召回；圖譜關係本身不等同於事實正確性的保證。


知識表示至少區分餐廳、文字片段、情境屬性、來源、站位與路線等實體。
餐廳與文字片段以提及關係連結，文字片段與情境屬性以支持關係連結，
每個支持關係皆附來源與核對狀態；相互衝突的證據並存而不直接覆寫。
餐廳與站位的鄰近關係僅提供交通候選，是否可搭乘仍由路線與時間分析判定。
此設計使語意相符、來源支持及交通可行性保有各自的判定依據。

\section{可行性判定與時間模型}
\label{sec:feasibility}

\subsection{三值判定}

令 \(\mathcal{H}(q)\) 為本次需求的必要限制索引集合，\(E\) 為可用證據；查核器 \(h_j(r,p,q,E)\in\{T,F,U\}\)。對每項必要限制，T 表示在可用證據與規則下通過，F 表示已知不符合，U 表示不足以判斷。任一項 F 則排除該餐廳與路線組合；全為 T 才進入已驗證排序；沒有 F 但有 U 則列入暫定區。某餐廳的一條路線不符合時，仍可檢查其他允許路線，不能立即排除整間餐廳。

{\small\singlespacing
\renewcommand{\arraystretch}{1.25}
\setlength{\tabcolsep}{3pt}
\begin{longtable}{|P{\dimexpr 0.1400\linewidth-2\tabcolsep-1.2500\arrayrulewidth\relax}|P{\dimexpr 0.2900\linewidth-2\tabcolsep-1.2500\arrayrulewidth\relax}|P{\dimexpr 0.2800\linewidth-2\tabcolsep-1.2500\arrayrulewidth\relax}|P{\dimexpr 0.2900\linewidth-2\tabcolsep-1.2500\arrayrulewidth\relax}|}
\caption{必要條件的三值判定}\label{table:three_valued}\\
\hline
\textbf{限制} & \textbf{通過 T} & \textbf{不通過 F} & \textbf{不確定 U} \\ \hline
\endfirsthead
\hline
\textbf{限制} & \textbf{通過 T} & \textbf{不通過 F} & \textbf{不確定 U} \\ \hline
\endhead
餐費預算 & 價格區間上界不超過同單位預算 & 區間下界已超過預算 & 價格不明、單位不明或區間跨預算 \\ \hline
營業時間 & 抵達與預計離店皆落在同一有效營業區間 & 抵達或用餐區間已知不符 & 缺時段、抵達或必要用餐時間 \\ \hline
步行上限 & 所有必要步行段符合使用者指定的上限定義 & 單段或總步行量已知超過所設定上限 & 缺所需步行路段、時間或上限定義 \\ \hline
上車班次 & 可在班次到站前完成步行與緩衝，且符合候車上限 & 已知允許的班次皆趕不上或服務不符 & ETA 過期、時間依據不足或班次無法對齊 \\ \hline
轉乘與次數 & 次數符合，前段到達加步行與緩衝仍接得上後段 & 已知超過次數或允許範圍內皆接不上 & 後段只有班距或時間無法確認 \\ \hline
\end{longtable}
}


三值的聚合採下列規則：
\[
H(r,p,q,E)=
\begin{cases}
F, & \exists j\in\mathcal{H}(q):h_j=F,\\
T, & \forall j\in\mathcal{H}(q):h_j=T,\\
U, & \text{其餘情形}.
\end{cases}
\]
可進入排序的路線集合定義為
\(\mathcal{P}_{T}(r,q)=\{p\in\mathcal{P}(r,o,t_{\mathrm{dep}}):H(r,p,q,E)=T\}\)。
若該集合為空，該餐廳不列入已驗證排序；存在未知但未被否定的路線時，仍可保留於暫定區。


最後點餐時間若已知，依該時間檢查；若未知，沿用只檢查用餐結束不晚於打烊的規則，並揭露未確認。一般情境不因最後點餐未提供而額外推翻既定規則；若使用者明確要求確認可點餐，則另作必要限制並判 U。

\subsection{上車與轉乘}

令 \(t_{\mathrm{ready}}\) 為到達上車站並完成緩衝的時間，\(t_{\mathrm{bus}}\) 為對齊後的預估到站或可用發車時間，\(W_{\max}\) 為候車上限，上車須符合式~\ref{eq:boarding}。

\begin{equation}\label{eq:boarding}
\begin{aligned}
t_{\mathrm{ready}} &= t_{\mathrm{dep}}+T_{\mathrm{walk\_to\_stop}}+b_{\mathrm{board}},\\
t_{\mathrm{ready}} &\leq t_{\mathrm{bus}}\leq t_{\mathrm{ready}}+W_{\max}.
\end{aligned}
\end{equation}

轉乘第 \(i\) 段時，令 \(a_i\) 為前段抵達，\(w_i\) 為轉乘步行，\(b_i\) 為緩衝，\(d_{i+1}\) 為後段可用發車時間，須符合式~\ref{eq:transfer}。

\begin{equation}\label{eq:transfer}
a_i+w_i+b_i\leq d_{i+1}.
\end{equation}

正常狀態與 EstimateTime 有值只是判定輸入，不能單獨當作可搭乘。尚未發車也不一律當作 F，須再看可用發車時間。路線、子路線、方向、站位與必要站序須對齊；時間為倒數值時，須有可確認的參考時刻，才能換算與使用者時間比較。

\subsection{抵達與用餐}

交通分析串接步行、候車、乘車、轉乘與最後步行，得到抵達時間，再加入用餐時長。若即時候車或轉乘更新，後續抵達及離店時間須重算；API 已包含的路段時間不能重複相加。

\begin{equation}\label{eq:timeline}
\begin{aligned}
t_{\mathrm{arr}}(r,p) &= t_{\mathrm{dep}}+T_{\mathrm{initial\_walk}}\\
 &\quad+\sum_i(T_{\mathrm{wait},i}+T_{\mathrm{ride},i})\\
 &\quad+\sum_iT_{\mathrm{transfer\_walk},i}+T_{\mathrm{final\_walk}},\\
t_{\mathrm{leave}}(r,p) &= t_{\mathrm{arr}}(r,p)+D_{\mathrm{meal}}.
\end{aligned}
\end{equation}

\(D_{\mathrm{meal}}\) 優先由使用者提供；若採 90–120 分鐘為預設假設，這是計畫值，須清楚顯示並可修改。營業查核使用完整用餐區間，跨日與午晚餐分段先轉為帶日期的區間。

\section{餐廳排序與交通便利性}
\label{sec:ranking}

交通便利性由步行、候車、乘車與轉乘負擔表示，與語意相關性分開計算。先排除必要條件不符的路線，為每間候選餐廳保留合適方案，再對餐廳排序，避免多條路線使同一間餐廳重複占據結果。

採用情境相符加分與交通負擔扣分的架構，組合分數與餐廳分數如式~\ref{eq:ranking}。\(p^*(r)\) 是餐廳 \(r\) 在可通過路線中的最佳方案；最終按 \(S_{\mathrm{restaurant}}\) 排餐廳，不輸出只有交通相近卻缺乏情境依據的店家。

\begin{equation}\label{eq:ranking}
\begin{aligned}
S(r,p\mid q) &= \alpha s_{\mathrm{sem}}+\beta s_{\mathrm{ctx}}
-\gamma c_{\mathrm{food}}-\delta u_{\mathrm{walk}}\\
&\quad-\varepsilon u_{\mathrm{wait}}
-\zeta u_{\mathrm{ride\_transfer}}-\eta u_{\mathrm{evid}},\\
p^*(r) &= \operatorname*{arg\,max}_{p\in\mathcal{P}_{T}(r,q)}S(r,p\mid q),\\
S_{\mathrm{restaurant}}(r\mid q)&=S(r,p^*(r)\mid q).
\end{aligned}
\end{equation}


其中 \(s_{\mathrm{sem}}\) 表示評論與需求的語意相關性，
\(s_{\mathrm{ctx}}\) 表示與指定情境相符的來源支持程度，
\(c_{\mathrm{food}}\) 為餐費成本，
\(u_{\mathrm{walk}}\) 與 \(u_{\mathrm{wait}}\) 分別為必要步行與候車負擔，
\(u_{\mathrm{ride\_transfer}}\) 為乘車時間與轉乘負擔，
\(u_{\mathrm{evid}}\) 為已足以查核之證據的品質成本。
成本項可使用預先固定上下界的截斷線性轉換：
\[
\mathcal{N}_{a,b}(x)
=\min\!\left(1,\max\!\left(0,\frac{x-a}{b-a}\right)\right),\qquad b>a.
\]
各項使用的 \(a,b\) 及權重在比較前固定，不依每次候選集合重新縮放。

各項用開發資料決定的固定尺度轉為 0–1，權重非負且總和為 1，在後續比較前固定。\(u_{\mathrm{ride\_transfer}}\) 將乘車時間與轉乘負擔分開計算後依固定規則合成，避免只用總時間掩蓋多次轉乘。尺度、權重與語意召回門檻須版本化，本文不預設最優參數。

餐費估計沿用價格區間中點，僅供顯示與評分；硬預算仍以區間查核，跨預算為 U。\(u_{\mathrm{evid}}\) 只比較仍足以查核的證據品質，必要證據缺漏者不得靠小幅扣分留在已驗證排序。輸出最多三間餐廳，同分以固定 ID 排序。

\section{核心演算法}
\label{sec:algorithm}


\renewcommand{\lstlistingname}{演算法}   % 標題顯示 Algorithm 而不是 Listing
\lstset{
  basicstyle=\ttfamily\footnotesize,
  columns=fullflexible, keepspaces=true,
  frame=single, rulecolor=\color{black!35}, framerule=0.6pt,
  backgroundcolor=\color{black!4},
  framexleftmargin=6pt, framexrightmargin=6pt,
  framextopmargin=5pt, framexbottommargin=5pt,
  aboveskip=1em, belowskip=1em,
  captionpos=t, breaklines=true, breakatwhitespace=false,
  keywordstyle=\bfseries\color{blue!55!black},
  morekeywords={if,for,each,in,return,continue,while,is,empty,Input,Output},
  literate={←}{{$\leftarrow$}}1,   % 讓 ← 在 XeLaTeX 也能正確顯示
}
\begin{lstlisting}[caption={Context-first restaurant recommendation with bus analysis},label={alg:recommend}]
Input : User query Q, origin O, departure time t_dep
Output: Restaurants R, bus plans P, evidence and unknown items D

01: intent ← PARSE_WITH_SCHEMA(Q, O, t_dep)
02: if required fields are missing: return ASK_USER(intent)
03: if units or values are invalid: return EXPLAIN_INPUT_ERROR
04: candidates ← RETRIEVE_CONTEXT_MATCHED_RESTAURANTS(intent)
05: candidates ← DEDUPLICATE_BY_RESTAURANT(candidates)
06: if candidates is empty: return EXPLAIN_NO_CONTEXT_MATCH
07: evidence ← RETRIEVE_CONTEXT_AND_DOMAIN_EVIDENCE(candidates)
08: C ← TOP_K(REMOVE_CONFIRMED_STATIC_FAILURES(candidates, evidence))
09: verified, provisional, records ← empty collections
10: for each restaurant r in C:
11:     plans, tool_status ← GET_BUS_PLANS(O, r, t_dep)
12:     if tool_status cannot establish usable plans:
13:         RECORD_FAILURE_OR_UNKNOWN(r, tool_status, records); continue
14:     for each plan p in plans:
15:         bus_data ← FETCH_AND_ALIGN_RELEVANT_BUS_DATA(p)
16:         timeline ← BUILD_UPDATED_TIMELINE(r, p, intent, bus_data)
17:         h ← CHECK_REQUIRED_CONSTRAINTS(r, p, timeline, evidence)
18:         record ← SAVE_DECISION(r, p, h, timeline, evidence, bus_data)
19:         records.append(record)
20:         if ANY_FALSE(h): continue
21:         if ANY_UNKNOWN(h): provisional.append(r, p, record); continue
22:         verified.append(r, p, SCORE_CONTEXT_AND_TRANSPORT(record))
23: best_plan_per_restaurant ← KEEP_BEST_VERIFIED_PLAN(verified)
24: R ← TOP_3_RESTAURANTS(STABLE_SORT(best_plan_per_restaurant))
25: if R is empty: return NO_VERIFIED_RESULT(provisional, records)
26: explanation ← GENERATE_FROM_RECORDS(R, records)
27: if explanation fails fact checks: explanation ← TEMPLATE(R, records)
28: return R, corresponding bus plans, explanation, provisional, records
\end{lstlisting}


路線查詢失敗與已知無服務分開記錄；未取得回應不直接等於證明無路線。所有候選完成後，才判斷是否沒有已驗證餐廳。使用者要放寬預算或步行限制時，須明確修改需求，系統不靜默放寬。

\section{決策流程}
\label{sec:flow}


% 架構與流程圖共用配色；在 figure 群組外載入。
\usetikzlibrary{arrows.meta,shapes.geometric}
\definecolor{archLLM}{HTML}{DCE9F8}
\definecolor{archLLMd}{HTML}{2F6DB5}
\definecolor{archPro}{HTML}{E6F2E6}
\definecolor{archProd}{HTML}{3C8D4F}
\definecolor{archKB}{HTML}{FFF3D6}
\definecolor{archKBd}{HTML}{C58A12}
\definecolor{archSrc}{HTML}{FDE9DC}
\definecolor{archSrcd}{HTML}{D9622B}
\definecolor{archGray}{HTML}{4A4A4A}

\begin{figure}[H]
\singlespacing
\centering
\resizebox{\textwidth}{!}{%
\begin{tikzpicture}[font=\small,>=Stealth,
  st/.style={draw=archProd,fill=archPro,line width=0.8pt,rounded corners=4pt,
    minimum width=4.1cm,minimum height=0.95cm,text width=3.85cm,align=center,inner sep=4pt},
  llm/.style={st,draw=archLLMd,fill=archLLM},
  dc/.style={diamond,aspect=2.6,draw=archKBd,fill=archKB,line width=0.8pt,
    text width=2.85cm,align=center,inner sep=1pt},
  tm/.style={st,draw=archGray,fill=white,rounded corners=7pt},
  bad/.style={st,draw=archSrcd,fill=archSrc},
  arr/.style={-{Stealth[length=2.3mm,width=1.8mm]},line width=0.9pt,archGray},
  lab/.style={font=\scriptsize\bfseries,fill=white,inner sep=1pt,text=archGray}]
\node[tm] (a) at (0,0) {1 輸入情境、位置與時間};
\node[llm] (b) at (0,-1.6) {2 AI 解析、程式檢查欄位};
\node[dc] (c) at (0,-3.4) {3 需求完整有效？};
\node[bad,text width=2.4cm,minimum width=2.6cm] (ask) at (-3.75,-1.6) {補問或說明錯誤};
\node[st] (d) at (0,-5.2) {4 評論語意召回、餐廳去重};
\node[dc] (e) at (0,-7) {5 有相符候選？};
\node[bad,text width=2.4cm,minimum width=2.6cm] (empty) at (-3.75,-7) {說明未找到相關餐廳};
\node[st] (f) at (0,-8.9) {6 取得餐廳證據與候選};
\node[st] (g) at (6,0) {7 選取下一餐廳或路線\\查詢公車方案與資料};
\node[dc] (h) at (6,-1.9) {8 資料足以分析？};
\node[st] (i) at (6,-3.8) {9 建立交通與用餐時間線};
\node[dc] (j) at (6,-5.7) {10 必要條件 T／F／U};
\node[bad,text width=3.0cm,minimum width=3.2cm] (unknown) at (11.3,-1.9) {記錄工具失敗或未知};
\node[bad,text width=3.0cm,minimum width=3.2cm] (reject) at (11.3,-5.7) {F：排除此路線\\U：暫定並保留原因};
\node[st] (k) at (6,-7.5) {11 計算情境與交通分數};
\node[dc] (l) at (6,-9.4) {12 還有候選或路線？};
\node[st] (m) at (6,-11.3) {13 每間餐廳保留最佳方案\\依情境與交通排序};
\node[dc] (n) at (6,-13.2) {14 有已驗證餐廳？};
\node[bad,text width=3.0cm,minimum width=3.2cm] (none) at (11.3,-13.2) {說明無結果原因\\與未知項};
\node[llm] (o) at (6,-15.2) {15 最多三間餐廳、評論依據\\交通比較與搭乘方案};
\draw[arr] (a)--(b);
\draw[arr] (b)--(c);
\draw[arr] (c.west)-|node[lab,pos=0.2,above]{否}(ask.south);
\draw[arr] (ask.north)|-(-1.5,-0.8)-|(b.north);
\draw[arr] (c)--node[lab,right]{是}(d);
\draw[arr] (d)--(e);
\draw[arr] (e.west)--node[lab,above]{否}(empty.east);
\draw[arr] (e)--node[lab,right]{是}(f);
\draw[arr] (f.east)--(2.9,-8.9)--(2.9,1.4)-|(g.north);
\draw[arr] (g)--(h);
\draw[arr] (h.east)--node[lab,above]{否}(unknown.west);
\draw[arr] (h)--node[lab,right]{是}(i);
\draw[arr] (i)--(j);
\draw[arr] (j.east)--node[lab,above]{有 F／U}(reject.west);
\draw[arr] (j)--node[lab,right]{全 T}(k);
\draw[arr] (k)--(l);
\draw[arr] (unknown.east)--(13.5,-1.9)--(13.5,-9.4)--(l.east);
\draw[arr] (reject.south)--(11.3,-9.4)--(l.east);
\draw[arr] (l.west)--node[lab,above]{是}(3.3,-9.4)--(3.3,0)--(g.west);
\draw[arr] (l)--node[lab,right]{否}(m);
\draw[arr] (m)--(n);
\draw[arr] (n.east)--node[lab,above]{否}(none.west);
\draw[arr] (n)--node[lab,right]{是}(o);
\node[anchor=north west,text width=4.5cm,font=\scriptsize,align=left] at (-1.95,-11.0) {
藍色：語言模型與說明\\
綠色：程式處理\\
黃色：條件判斷\\
橘色：失敗、排除或未知\\[4pt]
某一路線不通過，仍繼續處理其他候選。使用者可修改需求後重新輸入。};
\end{tikzpicture}%
}
\caption[決策流程圖]{決策流程圖：先召回情境相符的餐廳，再取得公車方案並逐項查核，最後對餐廳排序。未知與明確不符合分開記錄，所有候選處理後才判斷是否有可推薦餐廳。}
\label{fig:flow}
\end{figure}


流程先完成情境召回，再分析公車，最後排餐廳。某一條路線失敗不立即終止其他候選；未知資料也有獨立輸出，不混在已確認符合的結果。

\section{預期成果}
\label{sec:expected}

預期交付能理解用餐情境、召回相關餐廳、比較臺中公車交通並產生最終推薦的原型，輸出包含評論依據、價格、用餐時間、交通負擔、搭乘方案與未知項。另提供可回溯來源的候選餐廳、交通分析與查核紀錄。

預期系統能在情境相符的候選中，利用公車交通分析協助使用者取捨，並透過三值與時間查核呈現可判斷與未知的條件。以上為預期方向，尚非實驗結果。

\section{後續擴充}
\label{sec:extensions}

完成初版並確認資料可用後，可延伸到出發後監控、延誤時重排與替代餐廳提醒，以及用餐結束後的回程。這些功能需要額外的行程狀態、位置、通知與未來班次證據，不列入初版核心驗收。

若實作監控，須以同一可識別班次的絕對預估時刻比較延誤，上車後仍監控剩餘路段，資料失效時撤除受影響的通過聲明。替代餐廳不自動覆寫使用者已選的目的地。若處理回程，不能把起站末班發車直接當成使用者上車站的時間，也不能只檢查用餐時長區間的兩個端點就保證整段皆可行。這些均須另訂規格與評估後才納入。
